\documentclass[a4paper]{article}
\usepackage[pdftex]{graphicx}
\usepackage[utf8]{inputenc}
\usepackage{enumerate}
\usepackage{icomma}
\usepackage{amssymb}
\usepackage{tikz}
\usepackage{href-ul}
\hypersetup{
	colorlinks=true,
	linkcolor=blue,
	urlcolor=blue}
\usepackage{geometry}
\geometry{a4paper, top=15mm, left=15mm, right=15mm, bottom=15mm,
	headsep=10mm, footskip=12mm}

\begin{document}
	\begin{tikzpicture}(10,3)
		\draw[ultra thick](2,0) --(10,0) -- (10,1.3) --(2,1.3) -- (2,0);
		\draw[fill=black](2,0)-- (0,.6) -- (2,.6) -- (2,0);
		\node at (6,.6) {\large Lösungsblatt von \href{https://www.okuyakl.de}{www.okuyakl.de}};
	\end{tikzpicture}
	\vspace{0.5 cm}
	
\noindent{\bf Aufgabe 1.}\\

\renewcommand{\arraystretch}{3}
\begin{tabular}{p{3 cm}p{3 cm}p{3 cm}p{3 cm}p{3 cm}}
	a) ${\sqrt{2}\over2}$	& b)$\sqrt{3}$ & c)$\sqrt{v}$ & d) ${\sqrt{10}\over 2}$ & e) ${2\cdot \sqrt{14} \over 7}$\\
	f) ${\sqrt{2}\over2}$ & g) ${\sqrt{3}\over 2}$ & h)${1\over b+c}$ & i) ${x\sqrt{3}\over 3}$ & j)$\sqrt{a+b}$\\
	k) $2$ & l) ${sqrt{n}\over n}$\\
	\hline
\end{tabular}
\vspace{0.5 cm}

\noindent{\bf Aufgabe 2.}\\

\renewcommand{\arraystretch}{3}
\begin{tabular}{p{5 cm}p{5 cm}p{5 cm}}
	a)$\sqrt{3}+\sqrt{2}$ & b) ${\sqrt{7}-\sqrt{3} \over 10}$ & c) ${2\sqrt{2}+\sqrt{6}\over 2 }$ \\
    d) ${ 3\sqrt{10} - 4 \over 44}$ & e) ${3\sqrt{2} + 1 \over 17}$ & f)${3- \sqrt{3}\over 2}$\\
\hline
\end{tabular}
\vspace{0.5 cm}

\noindent{\bf Aufgabe 3.}\\

\renewcommand{\arraystretch}{3}
\begin{tabular}{p{2.5 cm}p{2.5 cm}p{2.5 cm}p{2.5 cm}p{2.5 cm}p{2.5 cm}}
	a) $\mathbb{D}=\mathbb{R}_0^+$ & b) $\mathbb{D}=[ 5;\infty [$ & c) $\mathbb{D}=[-2;\infty[$ & d) $\mathbb{D}=]-\infty;0]$ & e) $\mathbb{D}=]-\infty;0]$ & f) $\mathbb{D}=\mathbb{R}$ \\
	g) $\mathbb{D}=\mathbb{R}^+$ & h) $\mathbb{D}=\mathbb{R}\setminus \{0\}$ & i) $\mathbb{D}=\mathbb{R}$ & j) $\mathbb{D}=\mathbb{R}_0^+$ & k) $\mathbb{D}=\mathbb{R}_0^+$ & l) $\mathbb{D}=]2;\infty[$ \\
	\hline
\end{tabular}
\vspace{0.5 cm}

\noindent{\bf Aufgabe 4.}\\

\renewcommand{\arraystretch}{3}
\begin{tabular}{p{5 cm}p{5 cm}p{5 cm}}
	a) $\mathbb{D}=[-1;\infty[ \qquad \mathbb{L}=\{-1\} $ & b) $\mathbb{D}=[-2;\infty[ \qquad \mathbb{L}= \{2\}$ & c) $\mathbb{D}=\mathbb{R}_0^+ \qquad \mathbb{L}=\{3\} $  \\
	d) $\mathbb{D}=[4;\infty[ \qquad \mathbb{L}= \{ \}$ & e) $\mathbb{D}=[4;\infty[ \qquad \mathbb{L}=\{ \} $ & f) $\mathbb{D}=\mathbb{R}\,  \setminus \, ]-2 ; 2[ \qquad \mathbb{L}=\{ \}$\\
	\hline
\end{tabular}
\vspace{0.5 cm}

\noindent{\bf Aufgabe 9.}\\
Das große Quadrat hat die Seitenlänge $\sqrt{12,25}=3,5~{\rm LE}$; das kleine hat die Seitenlänge $\sqrt{6,25}=2,5~{\rm LE}$. Die Ecken, durch die die Gerade läuft, haben dann die Koordinaten $E_1(-3,5|3,5)$ und $E_2(-6|2,5)$. 
Damit ist die Steigung m:
$$ m = {3,5-2,5 \over -3,5+6}=0,4$$
Der	y-Achsenabschnitt t ist:
$$3,5 = 0,4\cdot(-3,5) + t \quad \Rightarrow t = 4,9$$
Funktionsgleichung:
$$g: \quad y=0,4x+4,9$$	

\begin{center}
	\includegraphics[width=7 cm]{../../viecher/endcomic.pdf}
	
	Hier geht es zurück zum \href{https://www.okuyakl.de/math/m9wuiL115/aa115.pdf}{Aufgabenblatt}
\end{center}


\end{document}